\chapter{1971:Gerhard Herzberg}

\label{chap:chemistry1971}

The Nobel Prize in Chemistry for the year 1971 was awarded to Gerhard Herzberg.

\textbf{Motivation:}

for his contributions to the knowledge of electronic structure and geometry of molecules, particularly free radicals

\section{Gerhard Herzberg}

\subsection*{Early Life and Education}

Gerhard Herzberg was born on 1904-12-25 in Hamburg, Germany.

\subsection*{Career and Major Research}

Herzberg studied at the Technische Hochschule Darmstadt, where he received his doctorate in 1928. After emigrating to Canada, he worked at the University of Saskatchewan and then at the National Research Council in Ottawa. He used molecular spectroscopy to determine the electronic structure and geometry of molecules, including short-lived free radicals.

\subsection*{Nobel Prize-winning Work}

The work for which Gerhard Herzberg was awarded the Nobel Prize in Chemistry in 1971 was described by the Nobel Committee as: "for his contributions to the knowledge of electronic structure and geometry of molecules, particularly free radicals".

Herzberg measured and analyzed the spectra of diatomic and polyatomic molecules and of radicals such as CH, OH, and CH2, establishing their structures and energy levels with high precision.

\subsection*{Significance and Impact}

Herzberg's spectroscopic data provided the experimental basis for understanding molecular structure and for identifying molecules in planetary atmospheres and interstellar space.

\subsection*{Later Career and Honors}

Herzberg remained at the National Research Council in Ottawa for the rest of his career. He died on 1999-03-03 in Ottawa, Canada.

\subsection*{Legacy}

Herzberg's books on molecular spectra and structure are classic references that trained generations of spectroscopists.

\subsection*{Modern Developments}

Herzberg's spectroscopic studies of free radicals and diatomic molecules established molecular spectroscopy and provided the foundational data for understanding chemical bonds. His work opened the modern field of molecular astrophysics, enabling astronomers to identify molecules in interstellar space, and his techniques are used to study reaction intermediates and atmospheric chemistry.

\subsection*{Selected Publications}

"Molecular Spectra and Molecular Structure, Vol. I: Spectra of Diatomic Molecules" (1939)

"The Spectra and Structures of Simple Free Radicals" (1971)

\clearpage

\section*{Chapter Summary}

The 1971 prize recognized Gerhard Herzberg for determining the electronic structure and geometry of molecules, especially free radicals, by spectroscopy. His measurements are fundamental to molecular physics and astronomy.

